\section{Contribuição dos Estudos para os Requisitos}

O Quadro~\ref{qua:estudos-requisitos} complementa a matriz de origem com uma leitura por estudo. Os vínculos são decisões de projeto motivadas por observações documentais e pelos fundamentos do capítulo; não significam que os artigos formularam ou validaram esses requisitos. Cada estudo aparece uma vez. A linha sem fonte revalidada registra a pendência, sem inventar contribuição substantiva.

\begingroup\small\setlength{\tabcolsep}{3pt}\def\LTcaptype{quadro}
\setlength{\LTleft}{0pt}
\setlength{\LTcapwidth}{\textwidth}\renewcommand{\arraystretch}{1.12}
\begin{longtable}{|p{2.8cm}|p{2.4cm}|p{4.3cm}|p{5.2cm}|}
\caption{Observações dos estudos e sua relação com os requisitos.}\label{qua:estudos-requisitos}\\\hline
\textbf{Estudo e fonte} & \textbf{Grupos R} & \textbf{Observação documental} & \textbf{Consequência para a especificação} \\ \hline
\endfirsthead
\multicolumn{4}{c}{\quadroname\ \thequadro\ -- Continuação}\\\hline
\textbf{Estudo e fonte} & \textbf{Grupos R} & \textbf{Observação documental} & \textbf{Consequência para a especificação} \\ \hline
\endhead
\hline\endfoot\hline\endlastfoot
\citeonline{Math2021_001} (páginas físicas 2–5) & R02, R09, R13, R15, R20, R21, R24 & Classes operacionais, seleção do conjunto PISA e relatos divergentes sobre teste. & Exigir alvo, população e partições identificados, com métricas por classe; não equiparar classes aos níveis oficiais. \\ \hline
\citeonline{Implementation2025_000} (páginas físicas 2–3 e 7–8) & R09, R16, R18–R21, R24 & Predição de nota posterior; validação 92\%, teste 75\%, revocação da classe baixa 43\%. & Documentar horizonte e reamostragem; distinguir validação, teste e erros por classe. \\ \hline
\citeonline{Multimodels2020_002} (registro histórico de acesso) & R06, R14 (decisão autoral) & Original não disponível para nova conferência. & Contribuição substantiva não revalidada; preservar identificação e pendência de recuperação, sem requisito técnico específico derivado. \\ \hline
\citeonline{Analysis2022_003} (páginas físicas 2–5) & R06, R09, R16, R18–R20, R24 & Classificações binária e múltipla; divergência entre tabela e narrativa. & Documentar disciplina, atributos, alvo e avaliação. Conservar inconsistências sem escolher por suposição um resultado correto. \\ \hline
\citeonline{Design2025_004} (páginas físicas 1–5) & R03, R08, R09, R15, R19, R24 & Trajetórias por conhecimentos e perfil, com relato insuficiente da avaliação. & Relacionar sugestão e contexto curricular; examinar comparadores e incerteza. Não adotar ganhos como metas do protótipo. \\ \hline
\citeonline{Identifying2017_006} (páginas físicas 4–5, 9 e 11) & R09, R15, R18–R21, R24 & Classe derivada de escore TIMSS e descrição invertida das dobras. & Explicitar transformação do alvo e partições. Distinguir problema de relato de erro de execução não demonstrado. \\ \hline
\citeonline{Innovative2023_005} (páginas físicas 9–12) & R06, R09, R15, R16, R24 & Trajetórias e recursos; 120 iniciais, 50+52 válidos e gl=80 no pré-teste não reconciliados. & Informar perdas, grupos e denominador em cada saída; impedir aprovação de comparação com contagens incompatíveis. \\ \hline
\citeonline{Computational2017_008} (páginas impressas 44–48 e 52–56) & R02, R07, R09, R15, R16, R24 & Modelos cognitivos comparados com registros humanos secundários; domínios e frações reutilizados. & Separar domínio, tutor, agente simulado e dado humano; não duplicar amostras. Simulação científica difere de teste sintético de engenharia. \\ \hline
\citeonline{Machine2024_009} (páginas físicas 4–7) & R01, R09, R15, R20, R24 & Notas e atividades de 50 estudantes, sem validação independente descrita. & Definir quando atributos estão disponíveis e exigir avaliação independente antes de usar previsão para apoio. \\ \hline
\citeonline{Enhancing2025_012} (páginas físicas 10–15) & R08, R09, R15, R24 & Animação e aprendizagem cooperativa combinadas; diferenças iniciais e escola por condição. & Registrar conteúdo, intervenção conjunta e grupos; restringir inferência causal e preservar interpretação docente. \\ \hline
\citeonline{He2025_6915} (páginas físicas 11–16 e 31–35) & R09, R15, R20, R24 & Efeito agregado predito e aplicação escolar; denominadores e professores por condição. & Distinguir publicação, turma e estudante; registrar agrupamento e potenciais confundidores. Não tratar efeito agregado como proficiência individual. \\ \hline
\citeonline{Villegas2025_6916} (páginas físicas 13–26) & R03, R09, R16, R21, R24, R25 & Retorno adaptativo em STEM; precisão, progresso, tempo e percepções com janelas distintas. & Definir disciplina e período por saída; separar precisão técnica e resultado educacional; submeter clareza ao exame docente. \\ \hline
\citeonline{Ozseven2026_6917} (páginas físicas 7 e 10–12) & R09, R18–R21, R24 & Predição de acerto por interação, partição por estudante e métricas de tarefas diferentes. & Definir classe positiva e unidade; separar regressão/classificação e evitar interpretar revocação de acerto como detecção de risco. \\ \hline
\citeonline{UniversityMathematics2026_6919} (páginas físicas 7–8 e 12–19; leitura parcial) & R08, R09, R15, R24, R25 & Integração de matemática e metrologia; questionários, entrevistas e grupos de anos distintos. & Vincular conteúdo curricular e contexto; separar percepção, desempenho e mecanismo de adaptação, sem exigir biometria ou Fourier no protótipo. \\ \hline
\citeonline{Echeveria2025_6920} (páginas físicas 3 e 6–12) & R02, R05, R09, R15, R16, R20, R21, R24, R25 & Categorias de mentalidade; itens usados no alvo também entre preditores. & Documentar rótulo e disponibilidade das variáveis, verificar sobreposição e perdas; separar mentalidade e domínio matemático. R05 é salvaguarda autoral, sem dano demonstrado. \\ \hline
\end{longtable}\endgroup
\fonte{Elaborado pelo autor (2026).}
